% ============================================================
% Section 57: 一维离子晶体的德拜温度与格林艾森常量
% ============================================================
\section{固体理论：一维离子晶体的德拜温度与格林艾森常量}
\label{sec:1d-ionic-debye-gruneisen}

\begin{modelbox}[题目：一维离子晶体的德拜温度与格林艾森常量]
质量为 $M_+$ 和 $M_-$ 的正负离子等间距（距离为 $a$）排列成一维晶格，近邻两离子间的相互作用势为
\begin{equation}
u(r)=-\frac{e^2}{r}+\frac{e^2b^{n-1}}{nr^n},
\end{equation}
式中 $e$ 为电子电荷，$n$ 及 $b$ 为常数。求：
\begin{enumerate}
    \item 采用德拜模型，求晶格德拜温度；
    \item 计算晶格的格林艾森常量。
\end{enumerate}
\end{modelbox}

\begin{tipbox}[考点快速分析]
\begin{itemize}
    \item \textbf{平衡条件}：$u'(a)=0$ 确定 $a=b$，化简势函数参数
    \item \textbf{恢复力常数}：$\beta=u''(a)$，一维离子晶体的简谐近似基础
    \item \textbf{长波声速}：$v_p=\omega/q$（$q\to0$），德拜模型的核心输入
    \item \textbf{一维德拜模型}：$D(\omega)=\text{const}$，总模式数 $=2N$（每原胞2原子，各1自由度）
    \item \textbf{格林艾森常量}：$\gamma=-\dfrac{a}{2}\dfrac{u'''(a)}{u''(a)}=\dfrac{n+4}{2}$，非简谐效应的度量
\end{itemize}
\end{tipbox}

\begin{derivationbox}[标准解题过程]
\textbf{Step 0: 恢复力常数与平衡条件}

\textit{相互作用势}：
\begin{equation}
u(r)=-\frac{e^2}{r}+\frac{e^2b^{n-1}}{nr^n}.
\end{equation}

\textit{平衡条件}（$r=a$ 处势能取极小，一阶导数为零）：
\begin{equation}
\frac{du}{dr}\bigg|_{r=a}=\frac{e^2}{a^2}-\frac{e^2b^{n-1}}{a^{n+1}}=0
\quad\Longrightarrow\quad a=b.
\end{equation}

\textit{恢复力常数}（二阶导数）：
\begin{equation}
\beta=\frac{d^2u}{dr^2}\bigg|_{r=a}=-\frac{2e^2}{a^3}+\frac{(n+1)e^2b^{n-1}}{a^{n+2}}
=(n-1)\frac{e^2}{a^3}=(n-1)e^2b^{-3}.
\end{equation}

\textbf{Step 1: 色散关系与长波声速}

一维双原子链的色散关系（参见 \ref{sec:1d-diatomic-heat-capacity}）为
\begin{equation}
\omega_{\pm}^2=\frac{\beta}{\mu}\Bigl[1\pm\sqrt{1-\frac{4\mu}{M}\sin^2(qa)}\Bigr],
\end{equation}
其中原胞总质量 $M=M_++M_-$，约化质量 $\mu$ 满足 $1/\mu=1/M_++1/M_-$。

\textit{长波极限}（$q\to0$）：声学支近似为
\begin{equation}
\omega_-^2\approx\frac{2\beta}{M}(qa)^2.
\end{equation}

代入 $\beta=(n-1)e^2/a^3$：
\begin{equation}
\omega_-\approx\sqrt{\frac{2(n-1)e^2}{Ma}}\,q\equiv v_pq,
\qquad v_p=\sqrt{\frac{2(n-1)e^2}{Ma}}.
\end{equation}

\textbf{Step 2: 德拜温度}

\textit{一维德拜模型}：模式密度为常数。链长 $L=2Na$（$N$ 个原胞），波矢空间态密度 $L/(2\pi)$，考虑 $\pm q$：
\begin{equation}
D(\omega)\,d\omega=2\cdot\frac{L}{2\pi}\,dq=\frac{2Na}{\pi v_p}\,d\omega.
\end{equation}

总模式数等于系统总自由度 $2N$：
\begin{equation}
\int_0^{\omega_D}D(\omega)\,d\omega=\frac{2Na}{\pi v_p}\,\omega_D=2N
\quad\Longrightarrow\quad
\omega_D=\frac{\pi v_p}{a}.
\end{equation}

代入 $v_p$：
\begin{equation}
\omega_D=\frac{\pi}{a}\sqrt{\frac{2(n-1)e^2}{Ma}}.
\end{equation}

\textit{德拜温度} $\Theta_D=\hbar\omega_D/k_B$：
\begin{equation}
\boxed{\Theta_D=\frac{\hbar\pi}{k_Ba}\sqrt{\frac{2(n-1)e^2}{Ma}}},
\qquad M=M_++M_-.
\end{equation}

\textbf{Step 3: 格林艾森常量}

\textbf{从标准定义到微观形式的过渡}：三维中的标准定义为（见~\S\ref{sec:gruneisen-thermal-expansion}）
\begin{equation}
\gamma=-\frac{\partial\ln\omega}{\partial\ln V}=-\frac{V}{\omega}\frac{\partial\omega}{\partial V}.
\end{equation}

对于一维链，``体积''就是链长 $L=2Na$（$N$ 为原胞数，固定不变），故 $\gamma=-\partial\ln\omega/\partial\ln L$。

\textit{关键：}$\omega$ 作为 $L$ 的函数到底是什么？长波声学支（式~(290)）为
\begin{equation}
\omega_- = \sqrt{\frac{2\beta}{M}}\,a|q|.
\end{equation}
周期边界条件给出允许的波矢 $q_m=2\pi m/L=\pi m/(Na)$（$m$ 为整数），代入得
\begin{equation}
\omega_- = \sqrt{\frac{2\beta}{M}}\,a\cdot\frac{\pi|m|}{Na}
= \frac{\pi|m|}{N}\sqrt{\frac{2\beta}{M}}.
\end{equation}
\textbf{晶格常数 $a$ 恰好消掉}！因此，对某一个固定的振动模式（$m$ 固定），频率只通过恢复力常数 $\beta$ 依赖于 $a$：
\begin{equation}
\omega\propto\sqrt{\beta},\qquad\text{即}\quad \ln\omega = \tfrac12\ln\beta + \text{const}.
\end{equation}

由此得到
\begin{equation}
\frac{\partial\ln\omega}{\partial\ln L}
=\frac{\partial\ln\omega}{\partial\ln a}
=\frac12\frac{d\ln\beta}{d\ln a}
=\frac{a}{2\beta}\frac{d\beta}{da},
\end{equation}
于是
\begin{equation}
\gamma=-\frac{\partial\ln\omega}{\partial\ln L}=-\frac{a}{2\beta}\frac{d\beta}{da}.
\end{equation}

最后代入 $\beta=u''(a)$，则 $d\beta/da=u'''(a)$，即得一维链的微观公式：
\begin{equation}
\boxed{\gamma=-\frac{a}{2}\frac{u'''(a)}{u''(a)}}.
\end{equation}
（式中 $1/2$ 因子的来源：$\omega\propto\sqrt{\beta}$，求对数导数时甩出一个 $1/2$。）

\textit{计算高阶导数}：
\begin{align}
u'(r)&=\frac{e^2}{r^2}-\frac{e^2b^{n-1}}{r^{n+1}},\\
u''(r)&=-\frac{2e^2}{r^3}+\frac{(n+1)e^2b^{n-1}}{r^{n+2}},\\
u'''(r)&=\frac{6e^2}{r^4}-\frac{(n+1)(n+2)e^2b^{n-1}}{r^{n+3}}.
\end{align}

在 $r=a=b$ 处：
\begin{align}
u''(a)&=-\frac{2e^2}{a^3}+\frac{(n+1)e^2}{a^3}=\frac{(n-1)e^2}{a^3}=\beta,\\
u'''(a)&=\frac{6e^2}{a^4}-\frac{(n+1)(n+2)e^2}{a^4}
=-\frac{(n+4)(n-1)e^2}{a^4}.
\end{align}

代入格林艾森公式：
\begin{equation}
\gamma=-\frac{a}{2}\cdot\frac{-(n+4)(n-1)e^2/a^4}{(n-1)e^2/a^3}
=\frac{a}{2}\cdot\frac{n+4}{a}
=\boxed{\frac{n+4}{2}}.
\end{equation}
\end{derivationbox}

\begin{formulabox}[答案汇总]
\begin{center}
\begin{tabular}{cl}
\toprule
问题 & 结果 \\
\midrule
(1) 德拜温度 & $\Theta_D=\dfrac{\hbar\pi}{k_Ba}\sqrt{\dfrac{2(n-1)e^2}{Ma}}$，$M=M_++M_-$ \\[12pt]
(2) 格林艾森常量 & $\gamma=\dfrac{n+4}{2}$ \\[10pt]
\bottomrule
\end{tabular}
\end{center}
\end{formulabox}

\begin{insightbox}[物理图像与概念深化]
\textbf{1. 为什么 $a=b$？——平衡条件的物理}

势函数第一项 $-e^2/r$ 是库仑吸引（长程），第二项 $+e^2b^{n-1}/(nr^n)$ 是短程排斥（泡利不相容导致的芯电子排斥）。平衡时两者抵消：
\begin{equation}
\underbrace{\frac{e^2}{a^2}}_{\text{吸引力}}=\underbrace{\frac{e^2b^{n-1}}{a^{n+1}}}_{\text{排斥力}}.
\end{equation}

当 $a<b$ 时排斥力占优，$a>b$ 时吸引力占优，$a=b$ 恰为平衡点。参数 $b$ 的物理意义是：若排斥势的幂次 $n$ 不同，平衡点 $a$ 不同，但 $b$ 作为``特征长度''始终满足 $a=b$。

\textbf{2. 恢复力常数 $\beta=(n-1)e^2/a^3$ 的图像}

$\beta$ 是两个竞争效应的净结果：
\begin{itemize}
    \item 库仑吸引贡献 $-2e^2/a^3$（使原子``软''，易位移）
    \item 短程排斥贡献 $+(n+1)e^2/a^3$（使原子``硬''，难位移）
    \item 净值：$(n-1)e^2/a^3$
\end{itemize}

$n$ 越大，排斥势越陡峭（越像硬球），$\beta$ 越大，晶格越``硬''，德拜温度越高。这与直觉一致：$n\to\infty$ 对应刚球模型，$\beta\to\infty$。

\textbf{3. 格林艾森常量与势函数形状}

$\gamma=\dfrac{n+4}{2}$ 直接由势函数的三阶导数决定，反映了势阱的\textbf{非对称性}：
\begin{itemize}
    \item $n=1$：$\gamma=2.5$。但 $n=1$ 时 $\beta=0$（恢复力常数消失），晶格不稳定，此极限无物理意义
    \item $n=2$：$\gamma=3$。$\beta=e^2/a^3>0$，晶格稳定
    \item $n=9$（NaCl 型离子晶体，Born-Mayer 势）：$\gamma=6.5$
    \item $n\to\infty$：$\gamma\to\infty$。刚球势的非简谐性极端强烈
\end{itemize}

典型离子晶体的实验值：NaCl $\gamma\sim1.5$，KCl $\gamma\sim1.4$。这些值远小于 $n=9$ 的理论值，原因是实际晶体中电子云屏蔽和远程库仑作用使有效势更``软''。

\textbf{4. 一维 vs 三维格林艾森常量}

本题中格林艾森常量的定义 $\gamma=-\dfrac{a}{2}\dfrac{u'''}{u''}$ 是\textbf{一维链的微观定义}。在三维中，由于不同偏振方向的声速不同，通常定义\textbf{模式格林艾森常量}：
\begin{equation}
\gamma_{\vec q,s}=-\frac{V}{\omega_{\vec q,s}}\frac{\partial\omega_{\vec q,s}}{\partial V},
\end{equation}
再对所有模式取平均得到\textbf{平均格林艾森常量} $\bar\gamma$。本题采用爱因斯坦近似（所有模式频率相同），故模式格林艾森常量即为平均格林艾森常量。

\textbf{5. 德拜温度与材料参数的标度关系}

从结果中提取标度关系：
\begin{equation}
\Theta_D\propto\frac{1}{a}\sqrt{\frac{e^2}{Ma}}\propto a^{-3/2}M^{-1/2}.
\end{equation}

\begin{itemize}
    \item 晶格常数 $a$ 越小，离子间距越近，恢复力越强，$\Theta_D$ 越高
    \item 离子质量 $M$ 越大，惯性越大，$\Theta_D$ 越低
    \item 电荷 $e$ 越大，库仑作用越强，$\Theta_D$ 越高
\end{itemize}

这与实验趋势一致：LiF（$a$ 小、$M$ 小）$\Theta_D\sim732\,\mathrm{K}$；RbI（$a$ 大、$M$ 大）$\Theta_D\sim103\,\mathrm{K}$。
\end{insightbox}

\begin{thinkbox}[考场速记：一维离子晶体参数问题的四步法]
遇到``求德拜温度和格林艾森常量''类问题：
\begin{enumerate}
    \item \textbf{写势函数}，用 $u'(a)=0$ 定平衡位置，化简参数
    \item \textbf{求恢复力常数} $\beta=u''(a)$，写色散关系，取长波极限得声速 $v_p$
    \item \textbf{德拜模型}：一维 $D(\omega)=\text{const}$，$\int_0^{\omega_D}D(\omega)d\omega=2N$，解 $\omega_D\to\Theta_D$
    \item \textbf{格林艾森常量}：$\gamma=-\dfrac{a}{2}\dfrac{u'''(a)}{u''(a)}$，代入势函数高阶导数
\end{enumerate}
\textit{核心记忆}：$\beta=(n-1)e^2/a^3$，$v_p=\sqrt{2\beta a^2/M}$，$\omega_D=\pi v_p/a$，$\gamma=(n+4)/2$。
\end{thinkbox}
